\chapter{Conclusion}
\label{ch:conclusion}

This thesis developed SNPbag, a transformer-based model for phenotype
prediction from genotype data, and compared it against LDpred2-auto on
simulated binary phenotypes. The evaluation used $100000$ synthetic
individuals from the HAPNEST project with $14000$ SNPs from
chromosome~22 and phenotypes simulated under the logistic disease model
of \citet{Majumdar2018}.

The pre-training phase demonstrated that transformer-based self-supervised
learning can capture meaningful genomic structure. Masked genotype
prediction accuracy improved consistently from $84.18\%$ to $90.19\%$ as
the training set grew from $1000$ to $100000$ individuals, with
training and validation losses remaining close throughout, indicating no
overfitting. The continued improvement at $100000$ individuals suggests
the model has not yet saturated.

Despite effective pre-training, neither method achieved strong phenotype
prediction. SNPbag obtained an AUC of $0.54$ and LDpred2-auto an AUC of
$0.51$, both only marginally above the random baseline of $0.5$. Several
factors explain this. SNPbag was fine-tuned for a single epoch with a
frozen encoder, meaning the representations could not adapt to the
phenotype prediction task and the MLP had to learn the entire mapping from
fixed features in one pass. LDpred2-auto was designed for genome-wide
summary statistics but operated on only $14000$ SNPs from a single
chromosome, capturing a small fraction of the genetic variation it would
normally leverage. More fundamentally, genotype-based prediction is bounded
by the heritability $h^2$ of the trait, so a substantial fraction of the
phenotypic variation is driven by environmental factors that no genotype
model can capture. The combination of limited genomic coverage, restricted
fine-tuning, and the inherent heritability ceiling makes the near-random
performance unsurprising.

The key conclusion is that the bottleneck is scale, not architecture. The
pre-training results show that the transformer learns useful genomic
representations, and SNPbag achieved a marginally higher AUC than
LDpred2-auto despite operating under severe constraints. With greater
computational resources, future work could extend the model to longer
sequences, include more training individuals, and fine-tune the full
encoder, each of which addresses one of the limitations identified above.